\documentclass[11pt]{article}
% Compile with: pdflatex outline.tex (run twice for cross-references/ToC)
%
% NOTE ON DOCUMENT CLASS: Automated Software Engineering (Springer) submissions ultimately require
% Springer Nature's official class (e.g. sn-jnl.cls, "sn-basic" / numbered-reference layout), which
% is not vendored in this repository. This outline uses the standard article class as a portable
% stand-in so it compiles anywhere; swap in Springer's class and \maketitle-equivalent frontmatter
% commands when the official template is downloaded from the journal's submission site, mirroring
% how docs/research/jss/outline.tex uses elsarticle for the Elsevier-targeted companion submission.

\usepackage{amsmath}
\usepackage{amssymb}
\usepackage{hyperref}
\usepackage{enumitem}
\usepackage[margin=1in]{geometry}

\title{Graph-Based Detection of Architectural Anti-Patterns and Prescriptive Refactoring in
Distributed Publish--Subscribe Systems}
\author{Author Name(s) Here}
\date{}

\begin{document}
\maketitle

\begin{itemize}[leftmargin=*]
\item \textbf{Target Journal:} Automated Software Engineering (AuSE) --- Springer
\item \textbf{Target Venue:} Special Issue ``Intelligent techniques for CI/CD, DevOps, software
  evolution, technical debt analysis, and refactoring recommendation.''
\item \textbf{Target Topic:} \emph{Technical debt analysis} (the anti-pattern catalog, \S4),
  \emph{refactoring recommendation} (the prescriptive operators, \S6), and \emph{CI/CD/DevOps} (the
  delta-aware quality gate, \S7) --- three of the special issue's five named areas; no AutoML/NAS/LLM
  contribution is claimed.
\end{itemize}

\noindent\emph{Provenance note: this outline folds the previously separate anti-pattern catalog and
its empirical validation (\texttt{docs/antipatterns.md}) into the closed-loop
prescriptive-refactoring pipeline (\texttt{docs/prescription.md}) under one narrative --- detect
named, severity-tiered architectural anti-patterns, prescribe targeted verified refactorings, and
gate CI/CD on newly introduced regressions. The manuscript body at \texttt{draft.md} has been
rewritten to match this outline's detect$\to$prescribe$\to$gate structure.}

\begin{abstract}
Distributed publish--subscribe middleware (ROS~2, DDS, MQTT) decouples producers and consumers, but
the resulting indirect dependency structure obscures how component failures cascade, and this
architectural technical debt accumulates invisibly to code-level static analysis (SCA) tools. Unlike
object-oriented design, which has mature anti-pattern catalogs, publish--subscribe architectures have
no equivalent catalog, and existing structural diagnostic frameworks typically operate open-loop,
ranking components by criticality without naming the responsible pathology or verifying a remedy. We
address both gaps with SaG-Prescribe, a graph-based framework that (1) detects twenty-one named,
severity-tiered publish--subscribe anti-patterns and bad smells, including Single Point of Failure,
God Component, Broker Saturation, Chatty Pair, and QoS Policy Mismatch, via topological signatures
over thirteen structural metrics and adaptive thresholds, validated against independent
failure-simulation ground truth (Spearman rank correlation of 0.876; F1 score of 0.923); and (2)
compiles the flagged components and patterns into a transformation policy of three targeted
graph-mutation operators, logical topic splitting, physical anti-affinity reallocation, and transport
QoS contract hardening, each re-verified by the same cascade simulator before being surfaced as a
recommendation. The pipeline is operationalized as a delta-aware CI/CD quality gate blocking only
newly introduced anti-patterns relative to a Git merge base. Across eight benchmark scenarios, the
generated prescriptions reduce the System Risk Index in every case, with relative reductions between
1.4\% and 15.9\% (Wilcoxon signed-rank p = 0.0156). Component-level failure-impact reductions average
only +4.61\% and regress in two scenarios, motivating a per-edit acceptance filter as future work. The
full generate-verify loop completes within CI/CD-compatible time budgets for most scenarios.
\end{abstract}

\noindent\textbf{Keywords:} architectural anti-patterns; technical debt analysis; refactoring
recommendation; publish--subscribe middleware; CI/CD quality gates; failure cascade simulation

% ---------------------------------------------------------------------------
% MANUSCRIPT OUTLINE
% ---------------------------------------------------------------------------

\section{Introduction}

\begin{itemize}[leftmargin=*]
\item \textbf{1.1 Context and motivation:} Pub-sub decoupling as the backbone of cyber-physical/
  microservice/IoT systems; decoupling obscures failure propagation and the resulting architectural
  technical debt.
\item \textbf{1.2 Two open gaps:} (a) the Architecture-Code Gap and the absence of a named, testable
  anti-pattern catalog for pub-sub topologies (SCA tools are topology-blind; OO/microservices catalogs
  don't transfer to broker-mediated communication); (b) the open-loop refactoring-recommendation gap
  --- even topology-aware diagnostics (including our own prior SaG framework) rank criticality
  without naming the fault or verifying its remediation.
\item \textbf{1.3 Proposed solution:} A unified detect$\to$prescribe$\to$gate pipeline ---
  twenty-one-pattern catalog with empirically validated detection rules, feeding a closed-loop
  mutation-and-resimulation engine, operationalized as a delta-aware CI/CD gate.
\item \textbf{1.4 Contributions:} Four explicit, numbered items: (1) the empirically-validated
  anti-pattern catalog and detection methodology; (2) the closed-loop prescriptive pipeline and its
  three operators, with an honestly-scoped mapping from operators to the five catalog patterns they
  directly automate (of twenty-one total; see \S6.3); (3) the CI/CD gate; (4) the multi-scenario
  evaluation, including the honest component-level result.
\item \textbf{1.5 Relationship to the Authors' Prior Work:} Positions this submission relative to the
  companion Software-as-a-Graph JSS submission and the RASSE 2025 conference paper on graph-based
  critical-component identification; clarifies that the catalog and prescriptive material here are
  original to this manuscript and not duplicated review.
\item \textbf{1.6 Organization:} Standard structural overview of the remaining sections.
\end{itemize}

\section{Background and Related Work}

\begin{itemize}[leftmargin=*]
\item \textbf{2.1 Publish--Subscribe Middleware Dependability:} Broker fault tolerance, Kafka, DDS
  QoS/latency literature.
\item \textbf{2.2 Anti-Pattern and Code-Smell Catalogs:} Object-oriented design (Fowler's refactoring
  catalog; Brown et al.'s architectural AntiPatterns; Suryanarayana et al.'s design smells) and
  microservices smells (Richardson; Taibi et al.), establishing the template this catalog follows ---
  named pattern, formal detection rule, refactoring strategy --- while operating on system
  \emph{topology} rather than code or REST call graphs.
\item \textbf{2.3 Refactoring Recommendation and Architectural Technical Debt:} Code-scope
  recommenders (smell-driven, learning-based, LLM-based) vs. this paper's topology scope and
  verify-before-recommend model. Two of three prior citation slots are now fillable from \S2.2's real
  references; the learning-based and LLM-based recommender slots remain pending, with no invented
  references used.
\item \textbf{2.4 Search-Based Software Engineering and Architecture Optimization:} Open-loop vs.
  this paper's closed-loop verification.
\item \textbf{2.5 Diagnostic Foundation (SaG):} The heterogeneous graph model and RMAV attribution
  this paper builds on; what is summarized vs. not repeated from the companion diagnostic paper.
\item \textbf{2.6 Structural Criticality Analysis:} Why classical single-metric centrality (degree,
  betweenness alone) degrades on typed pub-sub graphs relative to the composite $Q(v)$ score.
\end{itemize}

\section{System Model and Code-Quality-Augmented Quality Attribution}

\begin{itemize}[leftmargin=*]
\item \textbf{3.1 Heterogeneous Graph Formulation:} Five node types, six structural edge types.
\item \textbf{3.2 Derived \texttt{DEPENDS\_ON} Projection:} App$\leftrightarrow$App,
  App$\leftrightarrow$Broker, App$\leftrightarrow$Library, Broker$\leftrightarrow$Broker, and the four
  architectural layers (app, infra, mw, system).
\item \textbf{3.3 Code Quality Penalty (CQP):} The paper's explicit bridge from SCA metrics to
  architecture-level risk; feeds the Maintainability dimension below.
\item \textbf{3.4 Multi-Dimensional Quality Attribution (RMAV):} Reliability/Maintainability/
  Availability/Vulnerability, AHP-weighted composite $Q(v)$, adaptive box-plot criticality tiers. This
  section is the shared foundation consumed by both the detection catalog (\S4) and the prescriptive
  engine (\S6).
\end{itemize}

\section{A Catalog of Architectural Anti-Patterns for Publish--Subscribe Systems}

\emph{(New detection contribution.)}

\begin{itemize}[leftmargin=*]
\item \textbf{4.1 Anti-Patterns vs. Bad Smells:} Taxonomy and confidence distinction --- structural
  configurations with well-understood failure modes vs. heuristic surface symptoms requiring human
  judgment.
\item \textbf{4.2 Detection Methodology:} Thirteen structural metrics per component (PageRank,
  Reverse PageRank, betweenness, closeness, eigenvector centrality, in/out-degree, clustering
  coefficient, articulation point score, bridge ratio, QoS-weighted degree measures); rank-based vs.
  min-max normalization; adaptive box-plot thresholds (scale-invariant, distribution-aware,
  statistically grounded) in place of fixed global constants.
\item \textbf{4.3 Catalog Overview:} Twenty-one patterns organized into three severity tiers
  (critical, high, medium) across four RMAV dimensions; summary table of pattern, severity, primary
  RMAV dimension, and formal detection rule, with full specifications and remediation strategies
  given in the companion technical reference (\texttt{docs/antipatterns.md}), cited rather than
  reproduced in full.
\item \textbf{4.4 Representative Pattern Walkthroughs:} Five patterns selected to span severity
  tiers, RMAV dimensions, and detection technique diversity: Single Point of Failure
  (articulation-point-based, Availability), God Component (betweenness plus maintainability gate),
  Broker Overload / Hub-and-Spoke (the pub-sub instantiation of a classical topology anti-pattern,
  validated on a deliberately-encoded scenario), Chatty Pair (bidirectional edge-weight product,
  hidden coupling behind nominal decoupling), and QoS Mismatch (transport-contract-boundary detection
  unique to QoS-bearing middleware).
\item \textbf{4.5 Empirical Validation of Detection:} Validation methodology (simulated failure
  impact as ground truth); headline metrics (Spearman correlation of 0.876 overall, rising to 0.943 at
  large scale; F1 of 0.923, precision of 0.912, recall of 0.857, Top-5 overlap of 0.80);
  pattern-specific evidence (Single Point of Failure F1 above 0.95; Hub-and-Spoke scenario
  confirmation); baseline comparison against single-metric centrality; the eight-scenario validation
  suite, foregrounding Scenario~06 (sparse microservices mesh) as the precision stress test --- a
  well-structured topology should produce few or no findings --- and Scenario~07 as the scalability
  benchmark.
\end{itemize}

\section{Closed-Loop Optimization Objective}

\begin{itemize}[leftmargin=*]
\item \textbf{5.1 Formal Objective:} A minimization objective under a currently unconstrained
  modification budget:
  \begin{equation*}
  \min_{\Delta} \sum_{v \in V} I^*_{\Delta(G)}(v) \quad \text{subject to} \quad \mathrm{Cost}(\Delta)
  \le \mathcal{B}
  \end{equation*}
  Acceptance criterion $\Delta\mathrm{SRI}>0$ as implemented in \texttt{PrescribeService}; explicit
  statement that the candidate set for $\Delta$ is exactly the critical/high components and patterns
  identified in \S4; note on the deferred margin-based criterion.
\end{itemize}

\section{The SaG-Prescribe Prescriptive Pipeline}

\begin{itemize}[leftmargin=*]
\item \textbf{6.1 Hexagonal Core Abstraction:} \texttt{IGraphRepository}, \texttt{Neo4jRepository}
  vs. \texttt{MemoryRepository}.
\item \textbf{6.2 Pipeline Stages 1--7:} Diagnostic foundation and anti-pattern detection $\to$
  prescriptive generation $\to$ review interface; explicit hand-off from \S4's named findings to
  operator selection.
\item \textbf{6.3 Three Refactoring Operators:} Logical topic splitting, physical anti-affinity
  reallocation, and transport QoS contract hardening --- each formalized as a typed graph mutation
  rule, triggered by two independent signals: a generic RMAV criticality tier (any critical/high
  component can trigger any operator) and detected-problem name matching, the only channel that ties
  a mutation back to a specific catalog pattern. Following the name-matching channel, only five of
  the twenty-one patterns are directly wired to an operator --- Single Point of Failure to
  anti-affinity reallocation; God Component, Bottleneck Edge, Failure Hub, and Hub-and-Spoke to topic
  splitting. Notably, QoS Mismatch has no wiring to QoS hardening despite the conceptual overlap ---
  QoS hardening fires only from the generic criticality tier. The remaining sixteen patterns have no
  automated operator and stay advisory-only; this asymmetry between the twenty-one-pattern catalog
  and the narrower five-pattern automation footprint is disclosed explicitly rather than implied
  away.
\item \textbf{6.4 Closed-Loop Verification Procedure:} Export $\to$ mutate $\to$ sandbox $\to$
  resimulate $\to$ $\Delta$SRI; policies are accepted or rejected as a whole --- there is no per-edit
  filter (see \S5.1) --- with the independence guarantee against circular leakage.
\end{itemize}

\section{DevOps Integration and Delta-Aware CI/CD Gating}

\emph{(This section is the paper's most direct link to the special issue's CI/CD/DevOps bullet ---
kept prominent.)}

\begin{itemize}[leftmargin=*]
\item \textbf{7.1 Automated Code Review Architecture:} \texttt{detect\_antipatterns.py} surfacing
  \S4's twenty-one detectors as a blocking CI check.
\item \textbf{7.2 Delta-Aware Regression Semantics:} Merge-base diffing, waiver register for
  accepted legacy risk.
\item \textbf{7.3 Exit-Code Protocol:} 0 (pass), 1 (warn), 2 (block).
\end{itemize}

\section{Experimental Design}

\begin{itemize}[leftmargin=*]
\item \textbf{8.1 Research Questions:}
  \begin{description}[leftmargin=2.2cm,style=nextline]
    \item[RQ1 (Detection efficacy and precision)] Does the catalog correlate with ground-truth
      impact, and does it avoid over-flagging well-structured systems?
    \item[RQ2 (Prescriptive efficacy)] Does the closed-loop engine reduce the System Risk Index
      across heterogeneous scenarios, and are the reductions statistically significant?
    \item[RQ3 (Operator contributions)] How do the individual refactoring operators contribute to
      the observed improvements across topological regimes?
    \item[RQ4 (Component-level remediation efficacy)] Do system-level SRI improvements translate
      uniformly into component-level failure-impact reductions, or can the greedy, unconditional
      policy application backfire on individual components?
    \item[RQ5 (Computational overhead and CI/CD feasibility)] What is the wall-clock execution time
      of the detection gate and of the full prescriptive pipeline as system scale grows, and is
      either compatible with CI/CD budgets?
  \end{description}
\item \textbf{8.2 Scenario Suites:} Eight scenarios (01--08) for detection validation, including the
  deterministic ``Tiny Regression'' smoke-test fixture (\S4.5); the seven-scenario subset (01--07,
  autonomous vehicle $\to$ hyper-scale enterprise) used for prescriptive evaluation, explicitly noting
  and justifying the count mismatch --- the smoke-test fixture carries no domain-representative
  signal for prescriptive evaluation.
\item \textbf{8.3 Experimental Protocol:} Deterministic simulator, verified
  $\sigma_{\text{seed}}=0$ across 5 seeds.
\item \textbf{8.4 Metrics:} Detection: Spearman $\rho$, $F_1$/precision/recall, Top-$k$ overlap.
  Prescription: SRI formula, $\Delta$SRI, operator counts, component-level deltas.
\end{itemize}

\section{Results}

\begin{itemize}[leftmargin=*]
\item \textbf{9.1 Detection Efficacy and Precision (RQ1):} Headline validation metrics from \S4.5;
  Scenario~06 precision-stress confirmation (few/no findings on a well-structured topology);
  Scenario~07 scalability confirmation.
\item \textbf{9.2 Prescriptive Efficacy (RQ2):} SRI improves in all 7 scenarios (1.4\%--15.9\%),
  Wilcoxon $W=0.0$, $p=0.0156$.
\item \textbf{9.3 Operator Contributions (RQ3):} Anti-affinity reallocation most frequent overall;
  QoS upgrades concentrate in high-loss profiles; contribution counts cross-referenced against the
  five name-matched patterns and the broader generic-criticality trigger described in \S6.3.
\item \textbf{9.4 Remediation Efficacy at Component Boundaries (RQ4):} Honest negative result: mean
  $+4.61\%$ component-level improvement but regressions in Hub-and-Spoke ($-31.67\%$) and
  Hyper-Scale ($-25.36\%$), traced to greedy unconditional policy application; motivates the per-edit
  acceptance filter.
\item \textbf{9.5 Computational Overhead and CI/CD Feasibility (RQ5):} Detection gate sub-2s to
  $\sim$40s across scales; full generate--verify loop $<$65s for 6/7 scenarios, $\sim$10.8 min at
  hyper-scale.
\end{itemize}

\section{Discussion, Threats to Validity, and Conclusion}

\begin{itemize}[leftmargin=*]
\item \textbf{10.1 What Naming and Verifying Buys:} Over an unnamed criticality score --- \S9.1 and
  \S9.4 read together.
\item \textbf{10.2 Positioning in CI/CD and Technical-Debt Workflows:} Advisory prescriptions vs.
  blocking detection gate; the catalog as an architecture-review checklist (design review by
  checklist, as in aviation or surgery) for teams without full CI/CD automation.
\item \textbf{10.3 Threats to Validity:} Construct validity (detection and verification both defined
  relative to the same discrete-event simulator; mitigation via operators/patterns grounded in
  established dependability practice; the operator-to-pattern linkage relies on substring matching of
  human-readable detected-problem names rather than a dedicated pattern-ID field, which is brittle to
  pattern-name changes and covers only five of the twenty-one catalog patterns directly, per \S6.3);
  internal validity (greedy, unconditional whole-policy application with no per-edit filter;
  independence guarantee against circular leakage; stratified reporting discipline); external
  validity (synthetic parameterized topologies; catalog scope limited to the pub-sub communication
  paradigm; QoS-mismatch detection currently specified for DDS/ROS~2 and MQTT weight semantics);
  conclusion validity (small scenario sample of $n=7$ for the Wilcoxon test; single-run,
  verified-deterministic simulator configuration).
\item \textbf{10.4 Limitations and Future Work:} Outlines seven prioritized next steps: (1) per-edit
  acceptance filter, highest priority, directly motivated by \S9.4; (2) budget-constrained policy
  search; (3) stochastic cascade propagation; (4) operator ablation and learned policy ordering; (5)
  catalog extension to hybrid REST/event and non-DDS/MQTT QoS semantics; (6) real-system replication
  (ATM topology); (7) LLM-assisted PR generation, kept to a single sentence and not oversold.
\item \textbf{10.5 Conclusion:} Summary of the end-to-end detect$\to$prescribe$\to$gate pipeline and
  headline results.
\end{itemize}

\section*{References}

\begin{itemize}[leftmargin=*]
\item The merged reference list draws on the eleven previously verified prescriptive/SBSE/
  pub-sub citations plus the anti-pattern catalog's own real, verified references (Fowler 1999; Brown
  et al. 1998; Suryanarayana et al. 2014; Richardson 2018; Taibi et al. 2020; Baldwin \& Clark 2000;
  Lehman 1996; Martin 2003; Nygard 2018; Colbourn 1987; Saaty 1980 --- deduplicated against the
  existing AHP citation). One citation slot in \S2.3 (learning-based refactoring-opportunity mining)
  and one (LLM-based refactoring suggestion) remain genuinely open and are not filled with invented
  references.
\end{itemize}

\section*{Pre-Submission Checklist}

\begin{itemize}[leftmargin=*]
\item Populate the two remaining reference slots in \S2.3 with real learning-based and LLM-based
  refactoring-recommendation citations. The smell-driven and architectural-technical-debt slots are
  now filled from \texttt{docs/antipatterns.md}'s own verified bibliography (Fowler 1999;
  Suryanarayana et al. 2014; Brown et al. 1998) --- no references have been invented for the
  remaining two.
\item Literature pass: expand from the current $\sim$20 references (11 prior + $\sim$9 newly merged
  from the catalog, after deduplicating Saaty) toward the $\sim$30--45 references AuSE reviewers
  typically expect.
\item Confirm companion-paper status (the Software-as-a-Graph JSS submission and the RASSE 2025
  conference paper on graph-based critical-component identification) at submission time and disclose
  per AuSE's originality/overlap policy in the cover letter. Verify the RASSE 2025 paper's scope does
  not overlap the newly added catalog material in \S4 beyond citation-level reuse of the shared
  RMAV/criticality foundation.
\item Reconcile scenario counts explicitly in \S8.2 --- eight scenarios for detection validation vs.
  seven for prescriptive evaluation is intentional (per \texttt{docs/antipatterns.md} \S6.3's note)
  but must be stated plainly on first mention to avoid a reviewer inconsistency flag.
\item Do not expand the LLM future-work mention (\S10.4, item 7) beyond a single sentence --- the
  paper has no LLM/AutoML contribution, and overselling it risks a scope-mismatch desk rejection.
\item Fill in Funding, Competing Interests, and Data Availability declarations before submission.
\end{itemize}

\end{document}
